\section{A counterexample to a tropical substitution rule}
\label{sec:tropical-comparison}

This appendix is independent of the game algorithm. It examines the
minimization variable-selection rule printed in Truffet's
\texttt{arXiv:2603.26423v4}, dated August~29, 2026
\cite{Truffet2026SubstitutionV4}. That version claims a strongly
polynomial tropical optimization algorithm and a mean-payoff consequence
through reductions. The example below loses its optimum under the stated
rule; the conclusion is confined to that version and procedure.

In ordinary arithmetic, consider
\begin{equation}\label{eq:substitution-example}
 \min x_1\quad\text{subject to}\quad
 \max(x_1,x_2)\ge0,\qquad x_1\ge\max(x_2+1,-2).
\end{equation}
The point $(0,-1)$ is feasible. Conversely, a feasible point with
$x_1<0$ would have $x_2\le x_1-1<0$, contradicting the first
constraint. Thus the true minimum is exactly zero.

\paragraph{Homogenization and the input conditions.}
The cited procedure homogenizes the constraint rows as
\[
 \max(x_1,x_2)\ge h,\qquad
 x_1\ge\max(x_2+1,h-2),
\]
and the objective as $z\ge x_1$; dehomogenization sets $h=0$.
Equation~(5b) of the cited paper holds for both ordered pairs of
constraint rows. Each left side lacks an $h$ term, whereas the other
row's right side has one. Setting $x_1=x_2=0$ and increasing $h$
therefore rules out the forbidden functional domination by any finite
additive constant. Both variables also have lower or upper occurrences, as required by
Definition~4.2 and Equation~(34). Choose the search endpoints
$\mu=0$ and $\lambda=-4$. They include the finite optimum witness
$(x_1,x_2,h)=(0,-1,0)$ and satisfy Equation~(23b) for the only
mixed lower function, $\max(x_2+1,h-2)$, since
\[
                 1+\lambda=-3<-2+\mu=-2.
\]
The lower function $h$ has no competing $x$ term.

\paragraph{The forced first selection.}
Numerotation convention~4.1 assigns index $0$ to the objective row.
The lower and upper occurrences in Equations~(24) and~(32) and
Definition~4.1 are drawn only from constraint rows $1,\ldots,m$.
For this instance they are
\[
 \begin{array}{c|c|c}
 \text{variable}&\text{lower-bound rows}&\text{upper-bound rows}\\\hline
 x_1&1,2&\varnothing\\
 x_2&1&2
 \end{array}
\]
In particular, the objective row is not an additional upper occurrence
of $x_1$. The lower branch for $x_2$ is $x_2\ge h$ in row~1; its
upper bound is $x_2\le x_1-1$ in row~2. Equation~(53a) therefore
makes the dominating-variable set exactly $\{x_2\}$. Equation~(53b)
restricts eligible lower inequalities to that set, leaving the unique
pair $(1,2)$. Theorem~4.3, Equation~(47), and Section~4.2.3,
Case~2.2.2, prescribe saturation of this selected inequality:
\[
 x_2=h.
\]
There is no tie-breaking choice in this step.

\paragraph{The resulting loss of the optimum.}
After substitution the second row requires $x_1\ge h+1$. The next
eligible saturation consequently sets $x_1=h+1$. The stopping test in
Equation~(39) fails before each of these two substitutions. At both
stages the objective still depends on $x_1$, so the separate rule for
an objective-independent switch to maximization does not apply.
Normalizing to $h=0$ gives objective value $1$, whereas the original
minimum in Equation~\eqref{eq:substitution-example} is $0$.

The first substitution already removes every optimum: selecting a
lower-bound disjunct on the basis of an upper occurrence of its variable
does not preserve the optimum of these coupled constraints. This
calculation does not address a different variable-selection rule, a
repaired algorithm, or another version. In particular, it is not an
assessment of the distinct 2025 HAL/MSR max-atom proof discussed in the
introduction. None of the results of that work or of the cited tropical
procedure is used in the proof of Theorem~\ref{thm:main}.
